% =====================================================================
%  Fase de construcao do GL-GRASP: GLGRASP::construction()
%
%  Para colar no TG, copie a partir de \section{...} e garanta que o
%  preambulo tenha os pacotes listados abaixo.
% =====================================================================
\documentclass[12pt,a4paper]{article}

\usepackage[utf8]{inputenc}
\usepackage[T1]{fontenc}
\usepackage[brazilian]{babel}
\usepackage[margin=2.5cm]{geometry}
\usepackage{amsmath, amssymb}
\usepackage{booktabs}
\usepackage{tabularx}
\usepackage{listings}
\usepackage{xcolor}
\usepackage{tikz}
\usetikzlibrary{shapes.geometric, arrows.meta, positioning}

\lstdefinestyle{cpp}{
  language=C++,
  basicstyle=\ttfamily\footnotesize,
  keywordstyle=\bfseries,
  frame=single,
  rulecolor=\color{gray!60},
  backgroundcolor=\color{gray!5},
  columns=fullflexible,
  breaklines=true,
  showstringspaces=false,
  aboveskip=6pt,
  belowskip=6pt
}
\lstset{style=cpp}

\tikzset{
  passo/.style = {rectangle, draw=black, thick, text width=4.6cm,
                  minimum height=0.8cm, align=center, font=\small},
  func/.style  = {rectangle, draw=black, thick, fill=gray!10,
                  text width=5.6cm, minimum height=0.8cm, align=center,
                  font=\small\ttfamily},
  dec/.style   = {diamond, aspect=2.2, draw=black, thick, text width=2.2cm,
                  align=center, font=\small, inner sep=1pt},
  ini/.style   = {rectangle, rounded corners=8pt, draw=black, thick,
                  fill=gray!20, text width=2.4cm, minimum height=0.7cm,
                  align=center, font=\small\bfseries},
  seta/.style  = {-{Stealth[length=2.5mm]}, thick},
  rot/.style   = {font=\scriptsize}
}

\newenvironment{artigo}
  {\begin{quote}\small\itshape}
  {\end{quote}}

\begin{document}

\section{Fase de construção do GL-GRASP}
\label{sec:construction}

\subsection{Objetivo}

A função \texttt{GLGRASP::construction()} gera uma solução completa a
partir do desenho original: insere os vértices incrementais um a um,
escolhendo a cada passo \emph{qual} vértice inserir e \emph{em qual
posição}. Ela implementa a \textit{Construction Phase} da Figura~3 de
Charytitsch e Nascimento (2026), correspondente aos \textit{Steps} 1 a 7
da Seção~4.3, e é chamada pela função \texttt{run()} uma vez a cada
iteração do GRASP.

\subsection{Estruturas envolvidas}

\begin{table}[htbp]
\centering\small
\caption{Estruturas de dados da fase de construção e sua correspondência
com a notação do artigo.}
\label{tab:estruturas-construcao}
\begin{tabularx}{\textwidth}{@{}l X c@{}}
\toprule
Estrutura & Significado & Artigo \\
\midrule
\texttt{S} & solução parcial em construção & $I'$ \\
\texttt{CL} & vértices incrementais ainda não inseridos & $CL$ \\
\texttt{G[l][i]} & menor distância de \textit{embedding} entre o vértice
  e seus vizinhos já presentes em \texttt{S} & $\mathcal{G}(u)$ \\
\texttt{RCL} & subconjunto de \texttt{CL} com $\mathcal{G}$ abaixo do
  limiar & $RCL$ \\
\texttt{getalpha()} & parâmetro de gulosidade, sorteado em
  \texttt{algorithm()} & $\varphi$ \\
\bottomrule
\end{tabularx}
\end{table}

\subsection{Implementação}

\subsubsection{Inicialização}

\begin{lstlisting}
initialize_all_iteration_structure(I);
set_start_iteration(chrono::steady_clock::now());
MTRand rng(getseed());
\end{lstlisting}

Reinicia a lista de candidatos e os contadores da iteração anterior,
registra o instante inicial e cria o gerador aleatório da construção,
usado na extração de vértices da RCL e na escolha de posição.

\subsubsection{Steps 1 e 2 --- solução parcial inicial}

\begin{lstlisting}
allocate_solution(I, S);
\end{lstlisting}

Aloca \texttt{S}, insere cada vértice original na sua posição original
$\pi_0$ e marca \texttt{CL[l][i] = 0}. Os vértices incrementais
permanecem com \texttt{CL[l][i] = 1}, e suas posições em \texttt{S} ficam
vazias.

\begin{artigo}
Step 1: $I'$ contém apenas os vértices originais.
Step 2: a lista de candidatos inicial é $CL = \hat{V}$.
\end{artigo}

\subsubsection{Step 3 --- critério guloso}

\begin{lstlisting}
init_G(I, S);
\end{lstlisting}

Para cada vértice $u \in CL$, calcula
\begin{equation}
\mathcal{G}(u) = \min_{(u,v),\,(v,u)\,\in\,I'} d(v,u)
\label{eq:G}
\end{equation}
percorrendo as listas de adjacência \texttt{FWD} (vizinhos no nível
$\lambda+1$) e \texttt{BWD} (vizinhos no nível $\lambda-1$) e considerando
apenas vizinhos com \texttt{CL == 0}, isto é, já presentes na solução
parcial. Se $u$ não possui nenhum vizinho em \texttt{S}, define-se
$\mathcal{G}(u) = \infty$.

\subsubsection{Step 4 --- lista restrita de candidatos}

\begin{lstlisting}
rebuild_RCL(I, S);
\end{lstlisting}

Calcula o limiar sobre os valores finitos de $\mathcal{G}$ em $CL$,
\begin{equation}
\xi = \mathcal{G}_{min} + \varphi\,(\mathcal{G}_{max} - \mathcal{G}_{min}),
\label{eq:xi}
\end{equation}
e define $RCL = \{v \in CL : \mathcal{G}(v) \le \xi\}$. Vértices com
$\mathcal{G} = \infty$ só entram na RCL quando não há nenhum candidato
com valor finito.

\subsubsection{Laço principal}

\begin{lstlisting}
while (getTotalToAssign() > 0) {
    pair<unsigned, unsigned> p = extract_from_RCL(rng);
    unsigned new_pos = choose_position(I, S, p.first, p.second, rng);
    insert_node(I, S, p.first, p.second, new_pos);
    update_G_after_insertion(I, S, p.first, p.second);
    rebuild_RCL(I, S);
}
\end{lstlisting}

O laço executa enquanto $CL$ não estiver vazia, e cada iteração insere
exatamente um vértice, em quatro etapas.

\paragraph{Seleção do vértice.}
\texttt{extract\_from\_RCL} sorteia um vértice da RCL e devolve o par
\texttt{p = (nível, posição original)}.

\begin{artigo}
Step 5: selecione um vértice $v \in RCL$ aleatoriamente.
\end{artigo}

\paragraph{Seleção da posição.}
\texttt{choose\_position} constrói o conjunto $P$ de posições candidatas:
para cada nível adjacente que possua vizinhos de $v$ já em \texttt{S},
sorteia entre a posição do vizinho de menor distância e a do vizinho de
maior distância; em seguida acrescenta a posição média dos elementos de
$P$; por fim, escolhe um elemento de $P$ aleatoriamente. Cada posição
candidata é ajustada para a posição viável mais próxima, respeitando a
restrição de deslocamento máximo $k$.

\begin{artigo}
Step 5: (a) para cada camada vizinha, escolha aleatoriamente entre a
posição disponível mais próxima do vizinho mais próximo e a do vizinho
mais distante; (b) adicione a posição disponível mais próxima da média das
posições em $P$; (c) escolha arbitrariamente uma posição de $P$.
\end{artigo}

\paragraph{Inserção.}
\texttt{insert\_node} insere o vértice em \texttt{new\_pos}, deslocando o
bloco de vértices ocupado uma posição acima quando necessário, e define
\texttt{CL[l][i] = 0}.

\begin{artigo}
Step 6: remova $v$ de $CL$.
\end{artigo}

\paragraph{Atualização.}
O vértice inserido passa a ser um vizinho posicionado dos vértices
adjacentes que ainda estão em $CL$. Para cada um deles,
\texttt{update\_G\_after\_insertion} aplica
$\mathcal{G} \leftarrow \min(\mathcal{G}, d)$, onde $d$ é a distância do
arco que os liga ao vértice recém-inserido. Em seguida, a RCL é
reconstruída.

\begin{artigo}
Step 6: se $CL$ não estiver vazia, volte ao Step 3.
\end{artigo}

Como $\mathcal{G}$ é definido como um mínimo e a inserção acrescenta
exatamente um vizinho posicionado, a atualização é exata e custa
$O(\text{grau})$, em vez de recalcular $\mathcal{G}$ para todos os
vértices de $CL$. Essa é a otimização sugerida no próprio artigo.

\subsubsection{Verificação final}

\begin{lstlisting}
if (!S.is_feasible(I, getk()))
    cerr << " NOT FEASIBLE!!! " << endl;
\end{lstlisting}

Verifica se nenhum vértice original violou a restrição
$|\pi_0(v) - \pi(v)| \le k$. A verificação não foi acionada em nenhuma das
12\,150 execuções realizadas.

\begin{artigo}
Step 7: retorne $I'$.
\end{artigo}

\subsection{Exemplo}

Considere um grafo com dois níveis. No nível~0 estão os vértices
originais $a$ e $b$ e o incremental $e$; no nível~1, os originais $c$ e
$d$ e o incremental $f$. A solução parcial inicial é:

\begin{center}
\begin{tabular}{lccc}
\toprule
 & posição 0 & posição 1 & posição 2 \\
\midrule
nível 0 & $a$ & $b$ & --- \\
nível 1 & $c$ & $d$ & --- \\
\bottomrule
\end{tabular}
\end{center}

Os arcos e as distâncias de \textit{embedding} relevantes são
$d(e,d) = 0{,}40$, $d(b,f) = 0{,}05$ e $d(e,f) = 0{,}30$.

\paragraph{Critério inicial.} Considerando apenas vizinhos já presentes
em $I'$:

\begin{center}
\begin{tabular}{lll}
\toprule
Vértice & Vizinhos em $I'$ & $\mathcal{G}$ \\
\midrule
$e$ & $d$ (0,40); $f$ ainda não está em $I'$ & 0,40 \\
$f$ & $b$ (0,05); $e$ ainda não está em $I'$ & 0,05 \\
\bottomrule
\end{tabular}
\end{center}

\paragraph{RCL.} Com $\varphi = 0{,}5$, pela
Equação~\eqref{eq:xi}:
$\xi = 0{,}05 + 0{,}5\,(0{,}40 - 0{,}05) = 0{,}225$.
Logo, $RCL = \{f\}$.

\paragraph{Iteração 1.}
\begin{enumerate}
  \item Seleciona $f$.
  \item O vizinho de menor distância em $I'$ é $b$, com $\pi(b) = 1$; a
        posição-alvo é 1 no nível~1.
  \item Insere $f$ na posição 1; $d$ é deslocado para a posição 2, o que
        exige $k \ge 1$.
  \item Atualiza $\mathcal{G}$: $e$ é adjacente a $f$ com
        $d(e,f) = 0{,}30$ e $f$ agora pertence a $I'$, portanto
        $\mathcal{G}(e) = \min(0{,}40;\ 0{,}30) = 0{,}30$.
  \item $RCL = \{e\}$.
\end{enumerate}

\paragraph{Iteração 2.}
\begin{enumerate}
  \item Seleciona $e$.
  \item Seus vizinhos em $I'$ são $d$ ($\pi = 2$) e $f$ ($\pi = 1$); a
        posição é escolhida a partir deles.
  \item Insere $e$.
  \item $CL = \emptyset$; a construção termina.
\end{enumerate}

Na iteração~1, o valor de $\mathcal{G}(e)$ foi alterado pela inserção do
vértice anterior. Por isso a ordem de inserção afeta o resultado, e a
extração aleatória da RCL faz com que iterações distintas do GRASP
produzam soluções distintas.

\subsection{Fluxograma}

\begin{figure}[htbp]
\centering
\resizebox{0.78\textwidth}{!}{%
\begin{tikzpicture}[node distance=0.75cm]
  \node (ini)   [ini] {Início};
  \node (s12)   [func, below=of ini]   {allocate\_solution()};
  \node (s3)    [func, below=of s12]   {init\_G()};
  \node (s4)    [func, below=of s3]    {rebuild\_RCL()};
  \node (dec)   [dec,  below=1.0cm of s4] {$CL = \emptyset$?};
  \node (fim)   [ini,  right=1.6cm of dec] {Fim};
  \node (s5)    [func, below=1.0cm of dec] {extract\_from\_RCL()};
  \node (s5b)   [func, below=of s5]    {choose\_position()};
  \node (s6)    [func, below=of s5b]   {insert\_node()};
  \node (s6b)   [func, below=of s6]    {update\_G\_after\_insertion()};
  \node (s4b)   [func, below=of s6b]   {rebuild\_RCL()};

  \node [right=0.25cm of s12, rot, align=left] {Steps 1--2};
  \node [right=0.25cm of s3,  rot, align=left] {Step 3};
  \node [right=0.25cm of s4,  rot, align=left] {Step 4};
  \node [right=0.25cm of s5,  rot, align=left] {Step 5};
  \node [right=0.25cm of s5b, rot, align=left] {Step 5 (a)--(c)};
  \node [right=0.25cm of s6,  rot, align=left] {Step 6};
  \node [right=0.25cm of s6b, rot, align=left] {Step 6};
  \node [right=0.25cm of s4b, rot, align=left] {Step 4};

  \draw [seta] (ini) -- (s12);
  \draw [seta] (s12) -- (s3);
  \draw [seta] (s3)  -- (s4);
  \draw [seta] (s4)  -- (dec);
  \draw [seta] (dec) -- node[above, rot] {sim} (fim);
  \draw [seta] (dec) -- node[right, rot] {não} (s5);
  \draw [seta] (s5)  -- (s5b);
  \draw [seta] (s5b) -- (s6);
  \draw [seta] (s6)  -- (s6b);
  \draw [seta] (s6b) -- (s4b);
  \draw [seta] (s4b.west) -- ++(-1.1,0) |- (dec.west);
\end{tikzpicture}}
\caption{Fluxo de execução de \texttt{GLGRASP::construction()}, com a
correspondência de cada chamada aos \textit{Steps} da Seção~4.3 do artigo
de referência.}
\label{fig:construction}
\end{figure}

\end{document}
